\section{Conclusions}
\label{sec:conclusions}

We have provided a theoretical foundation for several empirical approaches commonly 
associated with the equivalent-layer technique in geophysics. Our analysis clarifies why 
the widely used method has been so successful in applications such as the joint processing 
of gravity-gradient components, the conversion between gravity disturbances and gravity-gradient 
tensors, and the use of dipole layers for magnetic data processing and interpretation.

Synthetic tests show that discrete equivalent layers estimated by inversion of potential-field 
data closely approximate the closed-form expressions of continuous planar layers, 
whose existence is guaranteed by the planar equivalent-layer theorem. These findings 
confirm that planar equivalent layers are not merely mathematical abstractions but carry 
tangible physical significance rooted in boundary-value problems of potential theory.

Finally, we illustrated the use of an alternative dipole-based equivalent-layer approach 
in the interpretation of an airborne magnetic survey over the Registro do Araguaia alkaline 
complex in the Goiás Alkaline Province, Brazil. The results demonstrate that this formulation 
can recover the horizontal position of magnetic sources without requiring prior assumptions 
about their total magnetization.